For a cuspidal eigenform $f$, let $\Pi_f$ denote its unitary
adelic normalization: at a good prime $q$ its Satake parameters
$\alpha_q,\beta_q$ satisfy
$a_q=\sqrt q(\alpha_q+\beta_q)$ and
$\alpha_q\beta_q=\chi(q)$.
Apply \cite[\S\S0.3--0.7, Theorem~A, pp.~409--411]{carayol}
to $\Pi_f^\vee$, and set $\rho_f=\sigma(\Pi_f^\vee)^\vee$.
Here $F=\Q$, the real weight is $2$, and the coefficient place
lies above $2$. The degree of $F$ is odd, so the additional
finite discrete-series-place hypothesis for even degree does
not apply. The local prime $p$ is odd.
Carayol's reciprocity sends geometric Frobenius to a uniformizer,
and his $\omega_p$ is $|\cdot|_p$.
His normalized-induction formula is
\[
 \sigma(\operatorname{Ind}(\xi_1,\xi_2))
 =\xi_1^{-1}\omega_p^{-1/2}\oplus
   \xi_2^{-1}\omega_p^{-1/2}.
\]
Thus if $\Pi_{f,p}=\operatorname{Ind}(\mu_1,\mu_2)$,
$\rho_f$ on arithmetic Frobenius has eigenvalues
$\sqrt p\,\mu_1(p),\sqrt p\,\mu_2(p)$ whenever these
characters are unramified. In particular its good polynomial
is $X^2-a_qX+q\chi(q)$. Both the automorphic contragredient and
the dual on the Galois side are part of this convention.

For a block with ramified character, its local conductor is
exactly one: it is at most the level exponent and at least the
conductor of its determinant.
Casselman's newvector theorem and conductor calculation
\cite[Theorem~1, pp.~302--306]{casselman} show that a special
representation has conductor one only for an unramified twist;
a ramified twist has conductor at least two. A supercuspidal
representation has conductor at least two. Consequently
$\Pi_{f,p}=\operatorname{Ind}(\mu_1,\mu_2)$ with $\mu_1$
unramified and $\mu_2$ of conductor one, after interchanging
the characters if necessary. Its newvector space is a line.
On inertia the ramified character is the prescribed quadratic
character.

